% --- Archon physics formalization source begin ---
% archon:physics
% archon:covers PhyXMiniProblems/problem_phyx_mini_0679.lean
% archon:source-report reports/phyx_mini/problem_phyx_mini_0679.source.json

\chapter{Physics problem phyx\_mini\_0679}
\label{ch:PhyXMiniProblems_problem_phyx_mini_0679}

\paragraph{Problem source.}
\#\# Physical scenario

You are standing on the ground at the origin of a coordinate system. An airplane flies over you with constant velocity parallel to the \textbackslash{}( x \textbackslash{}) axis and at a fixed height of \textbackslash{}( 7.60 \textbackslash{}times 10\textasciicircum{}3 \textbackslash{}text\{ m\} \textbackslash{}). At time \textbackslash{}( t = 0 \textbackslash{}), the airplane is directly above you so that the vector leading from you to it is \textbackslash{}( \textbackslash{}vec\{P\}\_0 = 7.60 \textbackslash{}times 10\textasciicircum{}3 \textbackslash{}hat\{\textbackslash{}mathbf\{j\}\} \textbackslash{}text\{ m\} \textbackslash{}). At \textbackslash{}( t = 30.0 \textbackslash{}text\{ s\} \textbackslash{}), the position vector leading from you to the airplane is \textbackslash{}( \textbackslash{}vec\{P\}\_\{30\} = (8.04 \textbackslash{}times 10\textasciicircum{}3 \textbackslash{}hat\{\textbackslash{}mathbf\{i\}\} + 7.60 \textbackslash{}times 10\textasciicircum{}3 \textbackslash{}hat\{\textbackslash{}mathbf\{j\}\}) \textbackslash{}text\{ m\} \textbackslash{}) as suggested in figure.

\#\# Question

Determine the magnitude and orientation of the airplane's position vector at \textbackslash{}( t = 45.0 \textbackslash{}text\{ s\} \textbackslash{}).

\#\# Answer choices

- A: A: \textbackslash{}( 1.34 \textbackslash{}times 10\textasciicircum{}4 \textbackslash{}text\{ m\} \textbackslash{})

- B: B: \textbackslash{}( 2.91 \textbackslash{}times 10\textasciicircum{}4 \textbackslash{}text\{ m\} \textbackslash{})

- C: C: \textbackslash{}( 1.43 \textbackslash{}times 10\textasciicircum{}4 \textbackslash{}text\{ m\} \textbackslash{})

- D: D: \textbackslash{}( 1.87 \textbackslash{}times 10\textasciicircum{}4 \textbackslash{}text\{ m\} \textbackslash{})

\#\# Figure caption (auxiliary; use the image as primary evidence)

The image illustrates a diagram involving a person on the ground, represented by a figure on the left. The person is looking up at an airplane flying at an angle. Two vectors are drawn to represent different lines of sight:

1. **Vector \textbackslash{}(\textbackslash{}vec\{P\}\_0\textbackslash{})**: This vertical arrow starts from the person and points directly upwards to a horizontal dashed line, indicating a reference direction perpendicular to the ground.

2. **Vector \textbackslash{}(\textbackslash{}vec\{P\}\_\{30\}\textbackslash{})**: This diagonal arrow extends from the person to the airplane, indicating the line of sight to the aircraft. The number "30" suggests the vector is at a 30-degree angle to the vertical vector or another specific reference.

The airplane is depicted flying along the dashed horizontal line above the person, indicating its flight path or altitude. The angle between the vectors suggests the person is observing the airplane at a specific angle upward.

\#\# Dataset metadata

Kinematics; Spatial Relation Reasoning, Predictive Reasoning

\paragraph{Recorded answer/context.}
C: C: \textbackslash{}( 1.43 \textbackslash{}times 10\textasciicircum{}4 \textbackslash{}text\{ m\} \textbackslash{})

\paragraph{Figure/image path.}
/root/proposal\_for\_physic/science-mango/phyx\_mini\_run/phyx\_data/test\_image/679.png

\paragraph{Formalization target.}
create a compiling Lean file with sorry bodies at `PhyXMiniProblems/problem\_phyx\_mini\_0679.lean`. The Lean declarations must preserve the physical quantities, dimensions or dimensional roles, figure labels, governing-law hypotheses, and final relation expressed by this problem.
Use Mathlib/Physlib names found through LeanExplore where available. If a physics API is missing, introduce faithful local abstractions rather than scalar placeholder aliases.

\begin{theorem}[Physics formalization target]
\label{thm:physics:phyx_mini_0679:target}
\lean{PhyXMiniProblems.ProblemPhyXMini0679.problem_phyx_mini_0679}
\uses{def:physics:phyx-mini-0679:phyxminiproblems-problemphyxmini0679-positioncomponentinmeters, def:physics:phyx-mini-0679:phyxminiproblems-problemphyxmini0679-positionmagnitudeinmeters, def:physics:phyx-mini-0679:phyxminiproblems-problemphyxmini0679-orientationabovepositivexaxisradians, def:physics:phyx-mini-0679:phyxminiproblems-problemphyxmini0679-orientationabovepositivexaxisdegrees, def:physics:phyx-mini-0679:phyxminiproblems-problemphyxmini0679-airplaneflyoversetup, def:physics:phyx-mini-0679:phyxminiproblems-problemphyxmini0679-matchesflyoverscenario, def:physics:phyx-mini-0679:phyxminiproblems-problemphyxmini0679-matchesproblemreadouts, def:physics:phyx-mini-0679:phyxminiproblems-problemphyxmini0679-matchessuppliedflyoverfigure, def:physics:phyx-mini-0679:phyxminiproblems-problemphyxmini0679-satisfiesuniformhorizontalflight, def:physics:phyx-mini-0679:phyxminiproblems-problemphyxmini0679-displayedmagnitudeinmeters, def:physics:phyx-mini-0679:phyxminiproblems-problemphyxmini0679-recordeddatasetanswer, def:physics:phyx-mini-0679:phyxminiproblems-problemphyxmini0679-roundstonearesthundredmeters, def:physics:phyx-mini-0679:phyxminiproblems-problemphyxmini0679-roundstonearesttenthdegree, def:physics:phyx-mini-0679:phyxminiproblems-problemphyxmini0679-isuniquenearestdisplayedmagnitude}
The assigned autoformalize agent should translate this physics problem into Lean declarations in the covered file, with theorem and lemma proof bodies written as `by sorry`.
\end{theorem}
\begin{proof}
This is an autoformalization task, not a proof task. Produce statements that can later be proved by the physics prover without weakening the physical model.
\end{proof}

% --- Archon declaration topology begin ---
\section{Declaration topology}
\begin{definition}[Planar Position Quantity]
  \label{def:physics:phyx-mini-0679:phyxminiproblems-problemphyxmini0679-planarpositionquantity}
  \lean{PhyXMiniProblems.ProblemPhyXMini0679.PlanarPositionQuantity}
  A signed planar position vector with physical dimension length
\end{definition}
\begin{proof}
  This is the declared model definition of Planar Position Quantity.
\end{proof}
\begin{definition}[Planar Velocity Quantity]
  \label{def:physics:phyx-mini-0679:phyxminiproblems-problemphyxmini0679-planarvelocityquantity}
  \lean{PhyXMiniProblems.ProblemPhyXMini0679.PlanarVelocityQuantity}
  A signed planar velocity vector with physical dimension length per time
\end{definition}
\begin{proof}
  This is the declared model definition of Planar Velocity Quantity.
\end{proof}
\begin{definition}[Time Quantity]
  \label{def:physics:phyx-mini-0679:phyxminiproblems-problemphyxmini0679-timequantity}
  \lean{PhyXMiniProblems.ProblemPhyXMini0679.TimeQuantity}
  A physical time coordinate relative to the stated t = 0 origin
\end{definition}
\begin{proof}
  This is the declared model definition of Time Quantity.
\end{proof}
\begin{definition}[Length Quantity]
  \label{def:physics:phyx-mini-0679:phyxminiproblems-problemphyxmini0679-lengthquantity}
  \lean{PhyXMiniProblems.ProblemPhyXMini0679.LengthQuantity}
  A nonnegative physical length, used for the fixed altitude
\end{definition}
\begin{proof}
  This is the declared model definition of Length Quantity.
\end{proof}
\begin{definition}[Diagram Axis]
  \label{def:physics:phyx-mini-0679:phyxminiproblems-problemphyxmini0679-diagramaxis}
  \lean{PhyXMiniProblems.ProblemPhyXMini0679.DiagramAxis}
  Cartesian axes of the observer-centered ground frame
\end{definition}
\begin{proof}
  This is the declared model definition of Diagram Axis.
\end{proof}
\begin{definition}[to Fin]
  \label{def:physics:phyx-mini-0679:phyxminiproblems-problemphyxmini0679-diagramaxis-tofin}
  \lean{PhyXMiniProblems.ProblemPhyXMini0679.DiagramAxis.toFin}
  \uses{def:physics:phyx-mini-0679:phyxminiproblems-problemphyxmini0679-diagramaxis}
  Coordinate index corresponding to a diagram axis
\end{definition}
\begin{proof}
  This is the declared model definition of to Fin.
\end{proof}
\begin{definition}[position Vector Readout]
  \label{def:physics:phyx-mini-0679:phyxminiproblems-problemphyxmini0679-positionvectorreadout}
  \lean{PhyXMiniProblems.ProblemPhyXMini0679.positionVectorReadout}
  \uses{def:physics:phyx-mini-0679:phyxminiproblems-problemphyxmini0679-planarpositionquantity}
  Read a planar position vector in a selected length unit
\end{definition}
\begin{proof}
  This is the declared model definition of position Vector Readout.
\end{proof}
\begin{definition}[position Component Readout]
  \label{def:physics:phyx-mini-0679:phyxminiproblems-problemphyxmini0679-positioncomponentreadout}
  \lean{PhyXMiniProblems.ProblemPhyXMini0679.positionComponentReadout}
  \uses{def:physics:phyx-mini-0679:phyxminiproblems-problemphyxmini0679-planarpositionquantity, def:physics:phyx-mini-0679:phyxminiproblems-problemphyxmini0679-diagramaxis, def:physics:phyx-mini-0679:phyxminiproblems-problemphyxmini0679-positionvectorreadout}
  Read one component of a planar position in a selected length unit
\end{definition}
\begin{proof}
  This is the declared model definition of position Component Readout.
\end{proof}
\begin{definition}[position Magnitude Readout]
  \label{def:physics:phyx-mini-0679:phyxminiproblems-problemphyxmini0679-positionmagnitudereadout}
  \lean{PhyXMiniProblems.ProblemPhyXMini0679.positionMagnitudeReadout}
  \uses{def:physics:phyx-mini-0679:phyxminiproblems-problemphyxmini0679-planarpositionquantity, def:physics:phyx-mini-0679:phyxminiproblems-problemphyxmini0679-positionvectorreadout}
  Read the Euclidean magnitude of a planar position in a selected unit
\end{definition}
\begin{proof}
  This is the declared model definition of position Magnitude Readout.
\end{proof}
\begin{definition}[velocity Component Readout]
  \label{def:physics:phyx-mini-0679:phyxminiproblems-problemphyxmini0679-velocitycomponentreadout}
  \lean{PhyXMiniProblems.ProblemPhyXMini0679.velocityComponentReadout}
  \uses{def:physics:phyx-mini-0679:phyxminiproblems-problemphyxmini0679-planarvelocityquantity, def:physics:phyx-mini-0679:phyxminiproblems-problemphyxmini0679-diagramaxis}
  Read one component of a planar velocity in selected length/time units
\end{definition}
\begin{proof}
  This is the declared model definition of velocity Component Readout.
\end{proof}
\begin{definition}[time Readout]
  \label{def:physics:phyx-mini-0679:phyxminiproblems-problemphyxmini0679-timereadout}
  \lean{PhyXMiniProblems.ProblemPhyXMini0679.timeReadout}
  \uses{def:physics:phyx-mini-0679:phyxminiproblems-problemphyxmini0679-timequantity}
  Read a physical time coordinate in a selected time unit
\end{definition}
\begin{proof}
  This is the declared model definition of time Readout.
\end{proof}
\begin{definition}[length Readout]
  \label{def:physics:phyx-mini-0679:phyxminiproblems-problemphyxmini0679-lengthreadout}
  \lean{PhyXMiniProblems.ProblemPhyXMini0679.lengthReadout}
  \uses{def:physics:phyx-mini-0679:phyxminiproblems-problemphyxmini0679-lengthquantity}
  Read a physical length in a selected length unit
\end{definition}
\begin{proof}
  This is the declared model definition of length Readout.
\end{proof}
\begin{definition}[position Component In Meters]
  \label{def:physics:phyx-mini-0679:phyxminiproblems-problemphyxmini0679-positioncomponentinmeters}
  \lean{PhyXMiniProblems.ProblemPhyXMini0679.positionComponentInMeters}
  \uses{def:physics:phyx-mini-0679:phyxminiproblems-problemphyxmini0679-planarpositionquantity, def:physics:phyx-mini-0679:phyxminiproblems-problemphyxmini0679-diagramaxis, def:physics:phyx-mini-0679:phyxminiproblems-problemphyxmini0679-positioncomponentreadout}
  Metre readout of one position component
\end{definition}
\begin{proof}
  This is the declared model definition of position Component In Meters.
\end{proof}
\begin{definition}[position Magnitude In Meters]
  \label{def:physics:phyx-mini-0679:phyxminiproblems-problemphyxmini0679-positionmagnitudeinmeters}
  \lean{PhyXMiniProblems.ProblemPhyXMini0679.positionMagnitudeInMeters}
  \uses{def:physics:phyx-mini-0679:phyxminiproblems-problemphyxmini0679-planarpositionquantity, def:physics:phyx-mini-0679:phyxminiproblems-problemphyxmini0679-positionmagnitudereadout}
  Metre readout of a position-vector magnitude
\end{definition}
\begin{proof}
  This is the declared model definition of position Magnitude In Meters.
\end{proof}
\begin{definition}[velocity Component In Meters Per Second]
  \label{def:physics:phyx-mini-0679:phyxminiproblems-problemphyxmini0679-velocitycomponentinmeterspersecond}
  \lean{PhyXMiniProblems.ProblemPhyXMini0679.velocityComponentInMetersPerSecond}
  \uses{def:physics:phyx-mini-0679:phyxminiproblems-problemphyxmini0679-planarvelocityquantity, def:physics:phyx-mini-0679:phyxminiproblems-problemphyxmini0679-diagramaxis, def:physics:phyx-mini-0679:phyxminiproblems-problemphyxmini0679-velocitycomponentreadout}
  Metre-per-second readout of one velocity component
\end{definition}
\begin{proof}
  This is the declared model definition of velocity Component In Meters Per Second.
\end{proof}
\begin{definition}[time In Seconds]
  \label{def:physics:phyx-mini-0679:phyxminiproblems-problemphyxmini0679-timeinseconds}
  \lean{PhyXMiniProblems.ProblemPhyXMini0679.timeInSeconds}
  \uses{def:physics:phyx-mini-0679:phyxminiproblems-problemphyxmini0679-timequantity, def:physics:phyx-mini-0679:phyxminiproblems-problemphyxmini0679-timereadout}
  Second readout of a physical time coordinate
\end{definition}
\begin{proof}
  This is the declared model definition of time In Seconds.
\end{proof}
\begin{definition}[length In Meters]
  \label{def:physics:phyx-mini-0679:phyxminiproblems-problemphyxmini0679-lengthinmeters}
  \lean{PhyXMiniProblems.ProblemPhyXMini0679.lengthInMeters}
  \uses{def:physics:phyx-mini-0679:phyxminiproblems-problemphyxmini0679-lengthquantity, def:physics:phyx-mini-0679:phyxminiproblems-problemphyxmini0679-lengthreadout}
  Metre readout of a physical length
\end{definition}
\begin{proof}
  This is the declared model definition of length In Meters.
\end{proof}
\begin{definition}[orientation Above Positive XAxis Radians]
  \label{def:physics:phyx-mini-0679:phyxminiproblems-problemphyxmini0679-orientationabovepositivexaxisradians}
  \lean{PhyXMiniProblems.ProblemPhyXMini0679.orientationAbovePositiveXAxisRadians}
  \uses{def:physics:phyx-mini-0679:phyxminiproblems-problemphyxmini0679-planarpositionquantity, def:physics:phyx-mini-0679:phyxminiproblems-problemphyxmini0679-positioncomponentinmeters}
  For a first-quadrant vector, its orientation above the positive x-axis is the arctangent of vertical component divided by horizontal component. This is a general geometric readout; it contains no problem-specific coordinate
\end{definition}
\begin{proof}
  This is the declared model definition of orientation Above Positive XAxis Radians.
\end{proof}
\begin{definition}[radians To Degrees]
  \label{def:physics:phyx-mini-0679:phyxminiproblems-problemphyxmini0679-radianstodegrees}
  \lean{PhyXMiniProblems.ProblemPhyXMini0679.radiansToDegrees}
  Convert a radian angle to its degree readout
\end{definition}
\begin{proof}
  This is the declared model definition of radians To Degrees.
\end{proof}
\begin{definition}[orientation Above Positive XAxis Degrees]
  \label{def:physics:phyx-mini-0679:phyxminiproblems-problemphyxmini0679-orientationabovepositivexaxisdegrees}
  \lean{PhyXMiniProblems.ProblemPhyXMini0679.orientationAbovePositiveXAxisDegrees}
  \uses{def:physics:phyx-mini-0679:phyxminiproblems-problemphyxmini0679-planarpositionquantity, def:physics:phyx-mini-0679:phyxminiproblems-problemphyxmini0679-orientationabovepositivexaxisradians, def:physics:phyx-mini-0679:phyxminiproblems-problemphyxmini0679-radianstodegrees}
  Orientation above the positive x-axis, expressed in degrees
\end{definition}
\begin{proof}
  This is the declared model definition of orientation Above Positive XAxis Degrees.
\end{proof}
\begin{definition}[Ground Observer Frame]
  \label{def:physics:phyx-mini-0679:phyxminiproblems-problemphyxmini0679-groundobserverframe}
  \lean{PhyXMiniProblems.ProblemPhyXMini0679.GroundObserverFrame}
  Qualitative specification of the observer-centered ground frame
\end{definition}
\begin{proof}
  This is the declared model definition of Ground Observer Frame.
\end{proof}
\begin{definition}[Figure Vector Label]
  \label{def:physics:phyx-mini-0679:phyxminiproblems-problemphyxmini0679-figurevectorlabel}
  \lean{PhyXMiniProblems.ProblemPhyXMini0679.FigureVectorLabel}
  Vector labels printed in the supplied raster
\end{definition}
\begin{proof}
  This is the declared model definition of Figure Vector Label.
\end{proof}
\begin{definition}[Figure Vector Direction]
  \label{def:physics:phyx-mini-0679:phyxminiproblems-problemphyxmini0679-figurevectordirection}
  \lean{PhyXMiniProblems.ProblemPhyXMini0679.FigureVectorDirection}
  Qualitative directions of the two arrows in the supplied raster
\end{definition}
\begin{proof}
  This is the declared model definition of Figure Vector Direction.
\end{proof}
\begin{definition}[Airplane Flyover Figure]
  \label{def:physics:phyx-mini-0679:phyxminiproblems-problemphyxmini0679-airplaneflyoverfigure}
  \lean{PhyXMiniProblems.ProblemPhyXMini0679.AirplaneFlyoverFigure}
  \uses{def:physics:phyx-mini-0679:phyxminiproblems-problemphyxmini0679-figurevectorlabel, def:physics:phyx-mini-0679:phyxminiproblems-problemphyxmini0679-figurevectordirection}
  Presentation-level information transcribed from the primary bitmap
\end{definition}
\begin{proof}
  This is the declared model definition of Airplane Flyover Figure.
\end{proof}
\begin{definition}[Airplane Flyover Setup]
  \label{def:physics:phyx-mini-0679:phyxminiproblems-problemphyxmini0679-airplaneflyoversetup}
  \lean{PhyXMiniProblems.ProblemPhyXMini0679.AirplaneFlyoverSetup}
  \uses{def:physics:phyx-mini-0679:phyxminiproblems-problemphyxmini0679-planarpositionquantity, def:physics:phyx-mini-0679:phyxminiproblems-problemphyxmini0679-planarvelocityquantity, def:physics:phyx-mini-0679:phyxminiproblems-problemphyxmini0679-timequantity, def:physics:phyx-mini-0679:phyxminiproblems-problemphyxmini0679-lengthquantity, def:physics:phyx-mini-0679:phyxminiproblems-problemphyxmini0679-diagramaxis, def:physics:phyx-mini-0679:phyxminiproblems-problemphyxmini0679-groundobserverframe, def:physics:phyx-mini-0679:phyxminiproblems-problemphyxmini0679-airplaneflyoverfigure}
  Independent physical data of the flyover. positionFromObserver is the airplane's position vector with tail at the stationary observer. Neither its value at the query time nor any answer choice is fixed by this structure
\end{definition}
\begin{proof}
  This is the declared model definition of Airplane Flyover Setup.
\end{proof}
\begin{definition}[Matches Flyover Scenario]
  \label{def:physics:phyx-mini-0679:phyxminiproblems-problemphyxmini0679-matchesflyoverscenario}
  \lean{PhyXMiniProblems.ProblemPhyXMini0679.MatchesFlyoverScenario}
  \uses{def:physics:phyx-mini-0679:phyxminiproblems-problemphyxmini0679-airplaneflyoversetup}
  Qualitative facts stated about the observer and the coordinate frame
\end{definition}
\begin{proof}
  This is the declared model definition of Matches Flyover Scenario.
\end{proof}
\begin{definition}[Matches Problem Readouts]
  \label{def:physics:phyx-mini-0679:phyxminiproblems-problemphyxmini0679-matchesproblemreadouts}
  \lean{PhyXMiniProblems.ProblemPhyXMini0679.MatchesProblemReadouts}
  \uses{def:physics:phyx-mini-0679:phyxminiproblems-problemphyxmini0679-positioncomponentinmeters, def:physics:phyx-mini-0679:phyxminiproblems-problemphyxmini0679-timeinseconds, def:physics:phyx-mini-0679:phyxminiproblems-problemphyxmini0679-lengthinmeters, def:physics:phyx-mini-0679:phyxminiproblems-problemphyxmini0679-airplaneflyoversetup}
  Numerical readouts stated in the problem. The query-time position, magnitude, orientation, and answer label do not occur here
\end{definition}
\begin{proof}
  This is the declared model definition of Matches Problem Readouts.
\end{proof}
\begin{definition}[Matches Supplied Flyover Figure]
  \label{def:physics:phyx-mini-0679:phyxminiproblems-problemphyxmini0679-matchessuppliedflyoverfigure}
  \lean{PhyXMiniProblems.ProblemPhyXMini0679.MatchesSuppliedFlyoverFigure}
  \uses{def:physics:phyx-mini-0679:phyxminiproblems-problemphyxmini0679-airplaneflyoverfigure}
  Raw qualitative evidence from image 679.png. No position at 45 s, vector magnitude, orientation, or answer choice is represented in this predicate
\end{definition}
\begin{proof}
  This is the declared model definition of Matches Supplied Flyover Figure.
\end{proof}
\begin{definition}[Satisfies Uniform Horizontal Flight]
  \label{def:physics:phyx-mini-0679:phyxminiproblems-problemphyxmini0679-satisfiesuniformhorizontalflight}
  \lean{PhyXMiniProblems.ProblemPhyXMini0679.SatisfiesUniformHorizontalFlight}
  \uses{def:physics:phyx-mini-0679:phyxminiproblems-problemphyxmini0679-timequantity, def:physics:phyx-mini-0679:phyxminiproblems-problemphyxmini0679-diagramaxis, def:physics:phyx-mini-0679:phyxminiproblems-problemphyxmini0679-positioncomponentreadout, def:physics:phyx-mini-0679:phyxminiproblems-problemphyxmini0679-velocitycomponentreadout, def:physics:phyx-mini-0679:phyxminiproblems-problemphyxmini0679-timereadout, def:physics:phyx-mini-0679:phyxminiproblems-problemphyxmini0679-lengthreadout, def:physics:phyx-mini-0679:phyxminiproblems-problemphyxmini0679-airplaneflyoversetup}
  Uniform straight-line kinematics, stated in arbitrary compatible length and time units. The vertical component of velocity vanishes and every position has the given altitude. These laws are general in both times and contain no special formula or numerical conclusion for 45 s
\end{definition}
\begin{proof}
  This is the declared model definition of Satisfies Uniform Horizontal Flight.
\end{proof}
\begin{lemma}[horizontal Velocity In Meters Per Second eq 268]
  \label{lem:physics:phyx-mini-0679:phyxminiproblems-problemphyxmini0679-horizontalvelocityinmeterspersecond-eq-268}
  \lean{PhyXMiniProblems.ProblemPhyXMini0679.horizontalVelocityInMetersPerSecond_eq_268}
  \uses{def:physics:phyx-mini-0679:phyxminiproblems-problemphyxmini0679-velocitycomponentinmeterspersecond, def:physics:phyx-mini-0679:phyxminiproblems-problemphyxmini0679-airplaneflyoversetup, def:physics:phyx-mini-0679:phyxminiproblems-problemphyxmini0679-matchesproblemreadouts, def:physics:phyx-mini-0679:phyxminiproblems-problemphyxmini0679-satisfiesuniformhorizontalflight}
  The two supplied position vectors determine a horizontal speed of 268 m/s
\end{lemma}
\begin{proof}
  The conclusion follows from the governing relations and figure readouts named in the statement of horizontal Velocity In Meters Per Second eq 268.
\end{proof}
\begin{lemma}[query Position Components]
  \label{lem:physics:phyx-mini-0679:phyxminiproblems-problemphyxmini0679-querypositioncomponents}
  \lean{PhyXMiniProblems.ProblemPhyXMini0679.queryPositionComponents}
  \uses{def:physics:phyx-mini-0679:phyxminiproblems-problemphyxmini0679-positioncomponentinmeters, def:physics:phyx-mini-0679:phyxminiproblems-problemphyxmini0679-airplaneflyoversetup, def:physics:phyx-mini-0679:phyxminiproblems-problemphyxmini0679-matchesproblemreadouts, def:physics:phyx-mini-0679:phyxminiproblems-problemphyxmini0679-satisfiesuniformhorizontalflight}
  Uniform motion for another 15 s places the airplane at horizontal coordinate 12060 m; fixed altitude leaves the vertical coordinate at 7600 m
\end{lemma}
\begin{proof}
  The conclusion follows from the governing relations and figure readouts named in the statement of query Position Components.
\end{proof}
\begin{definition}[Answer Choice]
  \label{def:physics:phyx-mini-0679:phyxminiproblems-problemphyxmini0679-answerchoice}
  \lean{PhyXMiniProblems.ProblemPhyXMini0679.AnswerChoice}
  Labels of the four displayed magnitude choices
\end{definition}
\begin{proof}
  This is the declared model definition of Answer Choice.
\end{proof}
\begin{definition}[displayed Magnitude In Meters]
  \label{def:physics:phyx-mini-0679:phyxminiproblems-problemphyxmini0679-displayedmagnitudeinmeters}
  \lean{PhyXMiniProblems.ProblemPhyXMini0679.displayedMagnitudeInMeters}
  \uses{def:physics:phyx-mini-0679:phyxminiproblems-problemphyxmini0679-answerchoice}
  Metre values printed beside the four answer labels
\end{definition}
\begin{proof}
  This is the declared model definition of displayed Magnitude In Meters.
\end{proof}
\begin{definition}[recorded Dataset Answer]
  \label{def:physics:phyx-mini-0679:phyxminiproblems-problemphyxmini0679-recordeddatasetanswer}
  \lean{PhyXMiniProblems.ProblemPhyXMini0679.recordedDatasetAnswer}
  \uses{def:physics:phyx-mini-0679:phyxminiproblems-problemphyxmini0679-answerchoice}
  Answer label recorded in the dataset metadata; this is not a premise
\end{definition}
\begin{proof}
  This is the declared model definition of recorded Dataset Answer.
\end{proof}
\begin{definition}[Rounds To Nearest Hundred Meters]
  \label{def:physics:phyx-mini-0679:phyxminiproblems-problemphyxmini0679-roundstonearesthundredmeters}
  \lean{PhyXMiniProblems.ProblemPhyXMini0679.RoundsToNearestHundredMeters}
  An actual magnitude rounds to a displayed hundred-metre value
\end{definition}
\begin{proof}
  This is the declared model definition of Rounds To Nearest Hundred Meters.
\end{proof}
\begin{definition}[Rounds To Nearest Tenth Degree]
  \label{def:physics:phyx-mini-0679:phyxminiproblems-problemphyxmini0679-roundstonearesttenthdegree}
  \lean{PhyXMiniProblems.ProblemPhyXMini0679.RoundsToNearestTenthDegree}
  An angle readout rounds to a displayed tenth of a degree
\end{definition}
\begin{proof}
  This is the declared model definition of Rounds To Nearest Tenth Degree.
\end{proof}
\begin{definition}[Is Nearest Displayed Magnitude]
  \label{def:physics:phyx-mini-0679:phyxminiproblems-problemphyxmini0679-isnearestdisplayedmagnitude}
  \lean{PhyXMiniProblems.ProblemPhyXMini0679.IsNearestDisplayedMagnitude}
  \uses{def:physics:phyx-mini-0679:phyxminiproblems-problemphyxmini0679-positionmagnitudeinmeters, def:physics:phyx-mini-0679:phyxminiproblems-problemphyxmini0679-airplaneflyoversetup, def:physics:phyx-mini-0679:phyxminiproblems-problemphyxmini0679-answerchoice, def:physics:phyx-mini-0679:phyxminiproblems-problemphyxmini0679-displayedmagnitudeinmeters}
  A displayed magnitude is at least as close as every alternative
\end{definition}
\begin{proof}
  This is the declared model definition of Is Nearest Displayed Magnitude.
\end{proof}
\begin{definition}[Is Unique Nearest Displayed Magnitude]
  \label{def:physics:phyx-mini-0679:phyxminiproblems-problemphyxmini0679-isuniquenearestdisplayedmagnitude}
  \lean{PhyXMiniProblems.ProblemPhyXMini0679.IsUniqueNearestDisplayedMagnitude}
  \uses{def:physics:phyx-mini-0679:phyxminiproblems-problemphyxmini0679-airplaneflyoversetup, def:physics:phyx-mini-0679:phyxminiproblems-problemphyxmini0679-answerchoice, def:physics:phyx-mini-0679:phyxminiproblems-problemphyxmini0679-isnearestdisplayedmagnitude}
  A displayed magnitude is the unique nearest answer choice
\end{definition}
\begin{proof}
  This is the declared model definition of Is Unique Nearest Displayed Magnitude.
\end{proof}
% --- Archon declaration topology end ---
% --- Archon physics formalization source end ---
